\documentclass{article}
\usepackage[utf8]{inputenc}
\usepackage{fontspec}
\providecommand{\aitexlanguage}[1]{}
\ifdefined\aitexfont
  \setmainfont[Ligatures=Common]{Noto Sans}
  \newfontfamily\arabicfont{Amiri}
  \newfontfamily\hebrewfont{Noto Sans Hebrew}
  \newfontfamily\devanagarifont{Noto Sans Devanagari}
  \newfontfamily\cjkfont{Noto Sans CJK SC}
  \newfontfamily\thaifont{Noto Sans Thai}
  \newfontfamily\laofont{Noto Sans Lao}
\else
  % XeTeX does not consult OSFONTDIR for name-based platform discovery. File
  % names make all three producers consume the exact same locked font bytes.
  \setmainfont[Path=fonts/,Ligatures=Common]{NotoSans-Latin-Test.ttf}
  \newfontfamily\arabicfont[Path=fonts/]{Amiri-Arabic-Test.ttf}
  \newfontfamily\hebrewfont[Path=fonts/]{NotoSansHebrew-Test.ttf}
  \newfontfamily\devanagarifont[Path=fonts/]{NotoSansDevanagari-Test.ttf}
  \newfontfamily\cjkfont[Path=fonts/]{NotoSansCJKsc-Test.otf}
  \newfontfamily\thaifont[Path=fonts/]{NotoSansThai-Test.ttf}
  \newfontfamily\laofont[Path=fonts/]{NotoSansLao-Test.ttf}
\fi

\begin{document}
\section*{A portable Unicode workload}
Ordinary UTF-8 and OpenType shaping: office affine.

Direct U+2223 marker: scalar∣boundary.

Direct U+00A0 marker: direct space.

Included TeX U+2223 marker: include∣boundary.

Included TeX U+00A0 marker: include space.

\input{included.bbl}

{\arabicfont\aitexlanguage{ar}السَّلَامُ عَلَيْكُمْ}

{\hebrewfont\aitexlanguage{he}שָׁלוֹם עוֹלָם}

{\devanagarifont\aitexlanguage{hi}नमस्ते भारत}

{\cjkfont\aitexlanguage{zh}日本語中文测试漢字かなカナ}

{\thaifont\aitexlanguage{th}ภาษาไทย ยินดีต้อนรับ}

{\laofont\aitexlanguage{lo}ສະບາຍດີ ພາສາລາວ}
\end{document}
